\documentclass[a4paper,10pt]{article}

\usepackage[czech]{babel}
\usepackage[T1]{fontenc}
\usepackage[utf8]{inputenc}
\usepackage{amsmath,amssymb}
\usepackage[margin=18mm]{geometry}
\usepackage{enumitem}
\setlength{\parindent}{0pt}
\setlength{\parskip}{4pt}
\setlist[enumerate]{leftmargin=*,itemsep=9pt,topsep=5pt}

\begin{document}
\section*{Exponenciální rovnice A/B -- řešení pro učitele}
K testu \texttt{03\_test\_ExponencialniRovnice\_AB.tex}, 2.~ročník SPŠ.
Všechny rovnice řešíme v $\mathbb{R}$, bez použití logaritmů.

\smallskip
\textbf{Původ úloh:}
{\footnotesize\ttfamily
ucebnice/02\_Funkce\_a\_rovnice/\\
Exponencialni\_funkce\_a\_rovnice/\\
Exponenciální rovnice.tex\par}

\begin{center}
\begin{tabular}{c|cc}
Úloha & Varianta A & Varianta B \\ \hline
1 & Úloha 1.1 g) & Úloha 1.1 q) \\
2 & Úloha 1.2 n) & Úloha 1.2 q) \\
3 & Úloha 1.4 c) & Úloha 1.4 t) \\
4 & Úloha 1.6 c) & Úloha 1.6 e)
\end{tabular}
\end{center}

\subsection*{Varianta A}
U všech čtyř rovnic platí $D=\mathbb{R}$.
\begin{enumerate}
\item $\displaystyle 3^x=\frac1{81}$.
Protože $81=3^4$, dostáváme
\[
3^x=3^{-4}\iff x=-4.
\]
Tedy $K=\{-4\}$. Zkouška: $3^{-4}=\frac1{81}$.

\item $\displaystyle 2\cdot2^x\cdot4^{2-x}=\frac8{2^{3x+1}}$.
Převedeme vše na mocniny o základu $2$:
\begin{align*}
2^{1+x+2(2-x)}&=2^{3-(3x+1)},\\
2^{5-x}&=2^{2-3x},\\
5-x&=2-3x,\\
2x&=-3 \iff x=-\frac32.
\end{align*}
Tedy $K=\{-\frac32\}$. Zkouška: pro $x=-\frac32$ mají obě strany hodnotu $2^{13/2}$.

\item $\displaystyle \frac{3^{x+4}}9+9\cdot3^x=\frac23$.
Vytkneme společnou mocninu:
\begin{align*}
3^x\left(\frac{3^4}{9}+9\right)&=\frac23,\\
18\cdot3^x&=\frac23,\\
3^x&=\frac1{27}=3^{-3}\iff x=-3.
\end{align*}
Tedy $K=\{-3\}$. Zkouška: $\frac39+9\cdot\frac1{27}=\frac13+\frac13=\frac23$.

\item $\displaystyle 9\cdot3^x+3^{-x}=10$.
Položíme $t=3^x>0$, takže $3^{-x}=\frac1t$:
\begin{align*}
9t+\frac1t&=10,\\
9t^2-10t+1&=0,\\
(9t-1)(t-1)&=0,\\
t&=\frac19\quad\text{nebo}\quad t=1.
\end{align*}
Obě hodnoty splňují $t>0$. Zpětně $3^x=3^{-2}$ nebo $3^x=3^0$, tedy
$K=\{-2;0\}$. Zkoušky: $9\cdot\frac19+9=10$ a $9+1=10$.
\end{enumerate}

\newpage
\section*{Exponenciální rovnice -- varianta B}
Řešení k testu \texttt{03\_test\_ExponencialniRovnice\_AB.tex}.
U všech čtyř rovnic platí $D=\mathbb{R}$.
\begin{enumerate}
\item $\displaystyle \left(\frac12\right)^x=4$.
Píšeme $4=(\frac12)^{-2}$, takže
\[
\left(\frac12\right)^x=\left(\frac12\right)^{-2}\iff x=-2.
\]
Tedy $K=\{-2\}$. Zkouška: $(\frac12)^{-2}=4$.

\item $\displaystyle \frac{3^x}{9^{x-2}}=\frac{27^x}{9\cdot3^{4-x}}$.
Převedeme všechny mocniny na základ $3$:
\begin{align*}
\frac{3^x}{3^{2x-4}}&=\frac{3^{3x}}{3^{2+4-x}},\\
3^{4-x}&=3^{4x-6},\\
4-x&=4x-6,\\
5x&=10 \iff x=2.
\end{align*}
Tedy $K=\{2\}$. Zkouška: $L(2)=\frac9{9^0}=9$,
$P(2)=\frac{27^2}{9\cdot3^2}=\frac{729}{81}=9$.

\item $\displaystyle 2^{x+1}+2^{x-1}+2^{x+3}=\frac{21}{8}$.
Vytkneme $2^{x-1}$:
\begin{align*}
2^{x-1}\left(2^2+1+2^4\right)&=\frac{21}{8},\\
21\cdot2^{x-1}&=\frac{21}{8},\\
2^{x-1}&=2^{-3},\\
x-1&=-3 \iff x=-2.
\end{align*}
Tedy $K=\{-2\}$. Zkouška: $\frac12+\frac18+2=\frac{21}{8}$.

\item $\displaystyle 2^{x+3}+2^{3-x}=20$.
Položíme $t=2^x>0$ a využijeme $2^{-x}=\frac1t$:
\begin{align*}
8t+\frac8t&=20,\\
2t^2-5t+2&=0,\\
(2t-1)(t-2)&=0,\\
t&=\frac12\quad\text{nebo}\quad t=2.
\end{align*}
Obě hodnoty jsou kladné. Z rovnic $2^x=2^{-1}$ nebo $2^x=2^1$ vychází
$K=\{-1;1\}$. Zkoušky: $4+16=20$ a $16+4=20$.
\end{enumerate}

\subsection*{Kontrola úplnosti řešení}
U všech rovnic mají základy mocnin kladnou hodnotu různou od $1$.
V úlohách 1--3 jsme použili ekvivalentní úpravy a prostotu exponenciální funkce.
V úloze 4 obou variant je $t>0$, a při násobení rovnic hodnotou $t$
nevznikl nepřípustný kořen $t=0$. Všechny získané kořeny byly ověřeny
v původních rovnicích. Varianty zkoušejí stejné čtyři postupy a mají
srovnatelnou časovou náročnost (přibližně 15--20 minut).
\end{document}
